\section{Two precise limits of the remaining positivity arguments}
\label{sec:s5-scaled-datum}

The following two results test proposed extensions of the comparison
argument.  The first is an actual nonzero $L$-datum.  The second is
a finite representation-ring certificate, without a claimed arithmetic
realization.  Neither supplies a smaller unresolved arithmetic gate
than MA-RB-002.

\paragraph{MA-FQ-007 (a genuine rescaled datum with abscissa below one).}
There is a nonzero $L$-datum $\mathcal F_2=(f_2,0,m_2)$ in the
exact sense of Booker, Definition~1.3, such that:
\begin{enumerate}
\item its full archimedean kernel is zero, its spectral
      multiplicities are integers, and its coefficients $f_2(N)$
      are nonnegative for every $N>1$ and vanish off prime powers;
\item its associated Dirichlet series has nonnegative coefficients
      and abscissa of convergence $3/4$;
\item its completed function has simple real poles at $1/4$ and
      $3/4$, and its spectral support is contained in
      $|\operatorname{Im}z|\le1/4$ but is not contained in
      $\mathbb R$;
\item it is \emph{not} positive in Booker's separate sense of
      having only finitely many negative spectral multiplicities.
\end{enumerate}
In particular, the real-height support assumption of MA-FQ-006
cannot be replaced by the assertion that the abscissa is less
than one, even while retaining all the other hypotheses of that
theorem.  This example is not an $S_5$ comparison datum: its
coefficients violate both the finite-image bound and the specified
odd-prime-power pattern of that comparison.

\begin{proof}
Let $\chi_5$ be the primitive even quadratic character modulo five.
Use the genuine difference datum already constructed in MA-LD-003:
\[
 \mathcal F_0=\mathcal F_\zeta-\mathcal F_{\chi_5}
             =(f_0,0,m_0),\qquad
 L_{\mathcal F_0}(w)=R_0(w)=\frac{\zeta(w)}{L(w,\chi_5)}.
 \tag{MA-FQ.60}
\]
The real gamma factors agree.  Booker's conductor term in
Examples~1.4 gives the exact value at one, while logarithmic
differentiation gives the remaining coefficients:
\[
 f_0(1)=\tfrac12\log5,\qquad
 f_0(n)=(1-\chi_5(n))\frac{\Lambda(n)}{\sqrt n}
                                      \quad(n>1).
 \tag{MA-FQ.61}
\]
Define
\[
 \begin{gathered}
 m_2(z)=m_0(2z),\qquad K_2=0,\\
 f_2(n^2)=2f_0(n)\quad(n\ge1),\qquad
 f_2(N)=0\quad\text{if $N$ is not a square}.
 \end{gathered}
 \tag{MA-FQ.62}
\]
We verify all four axioms, including the normalization at one.

For (A1), $f_2(1)=\log5$ is real.  For $N>1$,
\[
 \begin{gathered}
 |f_2(N)|\ll N^{-1/4}\log N,\\
 \sum_{N\le X}|f_2(N)|^2
     =4\sum_{n\le\sqrt X}|f_0(n)|^2
     \ll 1+(\log X)^3 .
 \end{gathered}
 \tag{MA-FQ.63}
\]
The first estimate implies every required bound
$f_2(N)\log^jN\ll_j1$.  For the second estimate use
$|f_0(n)|\le2(\log n)n^{-1/2}$ for $n>1$ and compare
the resulting harmonic sum with its integral.  Thus the
mean square is $O_\varepsilon(X^\varepsilon)$ for every
$\varepsilon>0$.  Axiom (A2) holds with the identically
zero kernel.

For (A3), the actual integral completed divisor of (MA-FQ.60)
is discrete and has the usual polynomial counting bound.
Scaling its ordinates by $1/2$ preserves both properties and
every integer multiplicity.  More explicitly, the total variation
of $m_2$ for $|\operatorname{Re}z|\le T$ equals that of $m_0$
for $|\operatorname{Re}z|\le2T$.  The old support lies in
$|\operatorname{Im}z|\le1/2$, including the two completed
endpoints, and the new support lies in
$|\operatorname{Im}z|\le1/4$.

For (A4), take its actual test $g\in C_c^\infty(\mathbb R)$
and $h(z)=\int_{\mathbb R}g(x)e^{ixz}\,dx$ with
$h(\mathbb R)\subset\mathbb R$.  Put
$\widetilde g(x)=2g(2x)$; its transform is
$\widetilde h(z)=h(z/2)$ and is again real on the real axis.
The original explicit formula, with $K_0=0$, gives
\[
 \begin{aligned}
 \sum_zm_2(z)h(z)
   &=\sum_zm_0(z)h(z/2)\\
   &=-2\operatorname{Re}\sum_{n\ge1}2f_0(n)g(2\log n)\\
   &=-2\operatorname{Re}\sum_{N\ge1}f_2(N)g(\log N).
 \end{aligned}
 \tag{MA-FQ.64}
\]
The sums include $n=N=1$.  Consequently the scaling also
has the correct conductor term, rather than only a formal
Euler identity.  There is no Jacobian in the multiplicities:
the first equality merely reindexes the discrete divisor.
This proves that (MA-FQ.62) is a genuine $L$-datum.

Booker's normalization of its associated function now gives,
initially in an absolute-convergence half-plane,
\[
 \begin{aligned}
 \log L_{\mathcal F_2}(s)
   &=\sum_{n\ge2}\frac{2f_0(n)}{\log n^2}
                              (n^2)^{1/2-s}\\
   &=\sum_{n\ge2}\frac{f_0(n)}{\log n}n^{1-2s}
     =\log R_0(2s-\tfrac12).
 \end{aligned}
 \tag{MA-FQ.65}
\]
It follows that the exact finite function and analytic conductor
are
\[
 L_{\mathcal F_2}(s)
      =\frac{\zeta(2s-1/2)}{L(2s-1/2,\chi_5)},\qquad
 Q_{\mathcal F_2}=e^{-2f_2(1)}=\frac1{25}.
 \tag{MA-FQ.66}
\]
An actual completed function with divisor $m_2$ is
\[
 \mathscr R_2(s)=5^{-s+1/4}
                    \frac{\zeta(2s-1/2)}{L(2s-1/2,\chi_5)}.
 \tag{MA-FQ.67}
\]
Indeed the original completed ratio is $5^{-w/2}R_0(w)$;
the common real gamma factor cancels exactly.  Under
$s\mapsto1-s$, its variable $w=2s-1/2$ changes to $1-w$,
so it has the original completed functional equation centered
at $s=1/2$.  At $s=1/2+iz$ its old height is $2z$, as
required by (MA-FQ.62).

For completeness, the nonzero value of $L(1,\chi_5)$ needed
at the endpoints can be checked without a new zero-location
input.  Its convergent, residue-class-grouped series has terms
\[
 \begin{aligned}
 L(1,\chi_5)
 &=\sum_{a\in\{1,6,11,\ldots\}}
       \left(\frac1a-\frac1{a+1}-\frac1{a+2}
                                      +\frac1{a+3}\right)\\
 &=\sum_{a\in\{1,6,11,\ldots\}}
       \left(\frac1{a(a+1)}-
                          \frac1{(a+2)(a+3)}\right)>0.
 \end{aligned}
 \tag{MA-FQ.68}
\]
The completed nontrivial Dirichlet function is therefore
nonzero at one and, by its functional equation, at zero.
The completed zeta function has a simple pole at each endpoint.
Thus the original completed ratio has simple poles at $w=0,1$,
and (MA-FQ.67) has simple poles at $s=1/4,3/4$.
Their heights are $z=i/4,-i/4$, respectively.  In particular
the real-height support hypothesis fails unconditionally.

The local Euler factors of $R_0$ are
\[
 \begin{cases}
  1,&\chi_5(p)=1,\\
  (1+p^{-w})/(1-p^{-w})
          =1+2\sum_{j\ge1}p^{-jw},&\chi_5(p)=-1,\\
  (1-5^{-w})^{-1},&p=5.
 \end{cases}
 \tag{MA-FQ.69}
\]
They have nonnegative coefficients.  If
$R_0(w)=\sum a_0(n)n^{-w}$, then
\[
 a_2(n^2)=a_0(n)n^{1/2}\ge0,
 \qquad a_2(N)=0\quad(N\text{ not a square}).
 \tag{MA-FQ.70}
\]
This series converges absolutely for $\operatorname{Re}s>3/4$.
It cannot have a smaller abscissa, because its sum has the
actual pole at $3/4$.  Nonnegativity identifies its ordinary
and absolute abscissae.  Hence its abscissa is exactly $3/4$.

By (MA-FQ.61)--(MA-FQ.62), all the coefficients $f_2(N)$
for $N>1$ are nonnegative and are supported on prime powers.
Nevertheless this is not a positive datum in Booker's
terminology.  MA-LD-003 proves, using the degree-zero case
of Booker's Theorem~1.7, that $m_0$ has infinitely many
negative and infinitely many positive values.  The bijective
scaling of the divisor preserves that property for $m_2$.
Thus Euler coefficient positivity has not been substituted
for the separate spectral finiteness hypothesis.

Finally, if $\chi_5(p)=-1$, then
\[
                    f_2(p^2)=\frac{4\log p}{\sqrt p}.
 \tag{MA-FQ.71}
\]
There are arbitrarily large such primes.  An elementary
argument suffices here: if their set were finite, choose
a positive integer $N$ with $N\equiv2\pmod5$ and
$N\equiv1\pmod p$ for each prime in that finite set.
The Chinese remainder theorem permits these congruences.
No prime in the list or equal to five divides $N$, but
$\chi_5(N)=-1$, contradicting the multiplicativity of
$\chi_5$ on the prime factorization of $N$.
The finite-image comparison has a uniform bound
$|f(p^j)|\le C(\log p)p^{-j/2}$ outside a finite set;
at $j=2$, (MA-FQ.71) violates it by an unbounded factor
$4\sqrt p/C$.  No finite choice of bad primes repairs
that defect.  Moreover $f_2(p)=0$ for every prime,
and, at $\chi_5(p)=-1$, the nonzero exponents are
$2,6,10,\ldots$, whereas the exceptional $S_5$ pattern
has nonzero coefficients at the odd exponents.  This
establishes the asserted distinction from an actual
comparison in the stipulated family.
\end{proof}

\paragraph{MA-RB-004 (a genuine-character certificate for all induced order inequalities).}
Let $G=S_5$ and retain its seven irreducible characters
$\mathbf1,\eta,V,V\eta,W,W\eta,X$, of respective dimensions
$1,1,4,4,5,5,6$.  For each integer $k\ge1$, the formal
off-center order assignment
\[
 \begin{gathered}
 o(V\eta)=o(W)=k,\\
 o(\mathbf1)=o(\eta)=o(V)=o(W\eta)=o(X)=0
 \end{gathered}
 \tag{MA-RB.40}
\]
has the genuine order character
\[
                    \Theta_\gamma=k(V\eta+W).
 \tag{MA-RB.41}
\]
For every subgroup $H\le G$ and every genuine character
$\chi$ of $H$, it satisfies
\[
 \begin{aligned}
 o(\operatorname{Ind}_H^G\chi)
   &=\langle\operatorname{Ind}_H^G\chi,\Theta_\gamma\rangle_G\\
   &=\langle\chi,\operatorname{Res}_H^G\Theta_\gamma\rangle_H
      \in\mathbb Z_{\ge0}.
 \end{aligned}
 \tag{MA-RB.42}
\]
Together with the formal automorphic numerator order $n=0$,
it attains equality in both $v_1+e\ge k$ and the regular
order bound $\operatorname{ord}_\gamma\zeta_L\ge9k$ of
MA-RB-002.  Thus adding only nonnegative induced-character
order constraints cannot strengthen either bound or exclude
this formal order pattern.

\begin{proof}
Extend $o$ additively to the representation ring.  The
irreducible characters are orthonormal, so (MA-RB.40) gives
$o(\psi)=\langle\psi,\Theta_\gamma\rangle_G$ for every
character $\psi$.  To verify the reciprocity equality directly,
use the finite induction formula and the fact that
$\Theta_\gamma$ is a class function:
\[
 \begin{aligned}
 \langle\operatorname{Ind}_H^G\chi,\Theta_\gamma\rangle_G
 &=\frac1{|G||H|}\sum_{x\in G}\sum_{h\in H}
                   \chi(h)\overline{\Theta_\gamma(xhx^{-1})}\\
 &=\frac1{|H|}\sum_{h\in H}
                   \chi(h)\overline{\Theta_\gamma(h)}
  =\langle\chi,\operatorname{Res}_H^G\Theta_\gamma\rangle_H .
 \end{aligned}
 \tag{MA-RB.43}
\]
Both characters in the last inner product are genuine.
Decomposing their finite-dimensional representations into
irreducibles makes that inner product a sum of products
of nonnegative integer multiplicities.  This proves
(MA-RB.42), without restricting $H$ to the previously
listed subgroups or $\chi$ to one-dimensional characters.

For the pointwise analytic order identities already proved
in MA-RB-002, the assignment is
\[
 (b,n,v_0,v_1,e,o(W),o(W\eta))=(0,0,0,k,0,k,0).
 \tag{MA-RB.44}
\]
It gives $o(A/L_W)=-k$, $o(B)=k$, and
$v_1+e=k$.  The previously checked symmetric-square
and Rankin--Selberg identities give
$q_\gamma=v_1+e-k=0$, $r_4=q_\gamma+n=0$, and
$r_5=r_4+v_0+o(W\eta)=0$.  Their holomorphy constraints
are therefore satisfied.  The regular-character order is
\[
 o(\operatorname{Reg}_G)
    =\Theta_\gamma(1)=4k+5k=9k.
 \tag{MA-RB.45}
\]
Equivalently, it is the regular Artin factorization with
irreducible multiplicities equal to their dimensions.

Under exact Artin induction, (MA-RB.42) is precisely the
formal order of $L_{L^H}(s,\chi)$ at the hypothetical
point.  In particular it satisfies every order
nonnegativity supplied by the holomorphy of an abelian
relative Hecke function, as well as every further
holomorphic induced genuine character.  Assign the same
orders at the distinct reflected point $1-\gamma$.
This is compatible with the order relation in the functional
equations; an off-center pair imposes no central parity
condition on the individual orders.  Consequently no
contradiction, strict improvement of either displayed bound,
or exclusion of the common-zero pattern follows from adding
only these nonnegative induced-order constraints.
\end{proof}

The certificate in MA-RB-004 does not construct a number
field, automorphic representation, or $L$-function realizing
these orders.  Additional analytic information, including
a suitable nonvanishing or a stronger independently proved
functorial constraint, could still rule them out.  The lemma
audits the scope of this induction/holomorphy argument; it
does not reduce the shared-real-zero gate of MA-RB-002.
Likewise the genuine $L$-datum of MA-FQ-007 is an exact
counterexample to the stated abstract extension, not to an
$S_5$ theorem or to the retained real-support criterion.

\paragraph{Primary locators and proof scope.}
The defining axioms, including the separate meaning of
\emph{positive}, are Booker, Definition~1.3, author p.~3;
the degree-one Artin/automorphic constructions and exact
conductor term are Examples~1.4(2)--(3), p.~4, in
\texttt{arXiv:1308.3067v2}.  Those two primary pages were
freshly read for the rescaling.  The group property of
$L$-data is stated on p.~4.  The earlier complete audit
of the degree-zero part of Theorem~1.7, pp.~15--17, is
MA-LD-003; that theorem is used here only for the final
spectral-positivity distinction, not to establish the four
axioms or the abscissa.  No new audit of the classical
degree-one continuation and functional equation is claimed.

The $S_5$ character notation, exact Artin induction, and
tensor identities in MA-RB-004 are the already checked
MA-S5-003/006 and MA-RB-002.  Formula (MA-RB.43) proves
the needed reciprocity calculation here.  The interpretation
as a genuine finite representation concerns nonnegative
character multiplicities; it supplies no arithmetic weight
or analytic nonvanishing assertion.

Booker, A. R. (2015). $L$-functions as distributions.
\emph{Mathematische Annalen, 363}(1--2), 423--454.
\url{https://doi.org/10.1007/s00208-015-1178-z}.
Author manuscript: \url{https://arxiv.org/abs/1308.3067v2}.
